% ============================================================
% Section 3: Methodology
% ============================================================

\section{Methodology}
\label{sec:method}

\subsection{Dataset}
\label{subsec:dataset}

We train and evaluate DocVerify on FantasyID~\citep{fantasyid2024}, a public
benchmark for ID-document forgery detection and localization from Idiap for the
DeepID 2025 Challenge~\citep{deepid2025challenge}. FantasyID comprises approximately
3\,284 labeled images of synthetic fantasy identity documents: \textit{bonafide}
(genuine) and \textit{attack} (tampered via face swap or text inpainting).

Each image is accompanied by JSON annotations specifying bounding-box regions
with \texttt{region\_provenance: altered}. We convert these to binary
segmentation masks at training resolution ($224 \times 224$ pixels). We
apply \texttt{GroupShuffleSplit} to all 3\,284 available FantasyID images
(grouping by document identity, $\text{test\_size}=0.15$,
\texttt{random\_state}=42) to obtain approximately 2\,791 development images
and 493 blind holdout images. \rev{This split was performed once,
before any hyperparameter search or model-design decision, and the holdout
was accessed only for the final blind evaluation.} The dataset is class-imbalanced:
2\,351 attack images (71.6\%) and 933 bonafide images (28.4\%); this
imbalance is preserved in all splits and not corrected via resampling or
loss weighting. \rev{The mild imbalance (2.5:1, attack-majority)
leaves PR-AUC robust; a \texttt{pos\_weight} sweep confirms no significant
effect.}

\subsection{Architecture}
\label{subsec:arch}

DocVerify is a multi-task encoder-decoder network for simultaneous
document-level classification and pixel-level forgery localization
(Fig.~\ref{fig:architecture}).

\subsubsection{Encoder}
The encoder is a six-block lightweight CNN inspired by \citet{patel2023dcnn},
a deepfake-detection architecture specifically designed for forgery artifact
extraction. A lightweight encoder is a deliberate design choice: the nested
CV protocol requires training ${\sim}2\,500$ models, so per-trial cost must
be kept low (${\sim}3.5$~min on our hardware). Heavier pretrained backbones
(e.g.\ ResNet-50, ViT) would make this scale infeasible; moreover, with only
${\sim}2\,800$ training images, large pretrained encoders risk importing
ImageNet biases of uncertain benefit here. The channel progression is $8 \to 16 \to 32 \to 64 \to 128 \to
256$ filters. All layers use LeakyReLU ($\alpha=0.2$)
and batch normalization. Blocks 1--4 use $3 \times 3$
convolutions; blocks 5--6 use $5 \times 5$ convolutions. Downsampling uses
average pooling after blocks 2 and~3 (to preserve boundary information) and
max pooling after blocks 4 and~5 (for translation invariance). Skip connections
are extracted from each block prior to pooling.

\subsubsection{Classification Head}
Global average pooling is applied to the $7 \times 7$ bottleneck, yielding a
256-dimensional vector passed through four fully-connected layers
(256$\to$32$\to$16$\to$16$\to$1) with dropout and LeakyReLU, producing a
classification logit $\hat{y}_{\text{cls}}$.

\subsubsection{Segmentation Head (U-Net Decoder)}
A five-stage U-Net decoder~\citep{unet2015} reconstructs the spatial mask via
bilinear upsampling and skip-connection concatenation. The channel width halves
at each stage from $C_{\text{dec}}$ (a tunable HPO parameter) down to
$C_{\text{dec}}/16$. A final $1 \times 1$ convolution produces the
segmentation logit $\hat{Y}_{\text{mask}}$.

\subsection{Loss Functions}
\label{subsec:loss}

The combined training loss is:
\begin{equation}
  \mathcal{L} = \mathcal{L}_{\text{cls}} + \lambda_{\text{mask}} \,
                \mathcal{L}_{\text{seg}},
  \label{eq:total_loss}
\end{equation}
where $\lambda_{\text{mask}} \in [0.5, 3.0]$ is included in the HPO search
space. $\mathcal{L}_{\text{cls}} = \text{BCEWithLogits}(\hat{y}_{\text{cls}},
y_{\text{cls}})$. The segmentation loss combines BCE and Dice:
\begin{equation}
  \mathcal{L}_{\text{seg}} = \text{BCE}(\hat{Y}, Y) + 1 -
    \frac{2 \sum_i p_i g_i + \epsilon}{\sum_i p_i + \sum_i g_i + \epsilon},
\end{equation}
as introduced for volumetric segmentation~\citep{milletari2016vnet} to address
foreground/background imbalance in small forgery regions.

All models are trained with AdamW and ReduceLROnPlateau
(factor $0.5$, patience~3). Mixed-precision training (FP16/FP32 via PyTorch AMP)
is used throughout.

\begin{figure*}[!t]
  \centering
  \includegraphics[width=\textwidth]{fig1_pareto_front}
  \caption{Validation objective space across all 500 HPO trials (10 folds
    $\times$ 50 MOTPE trials). Grey: evaluated; blue: Pareto-optimal;
    dashed: Pareto frontier. Red diamond: ideal point $(1,1)$; annotated point:
    selected configuration (distance $d=0.316$).}
  \label{fig:pareto}
\end{figure*}

\subsection{Multi-Objective HPO}
\label{subsec:hpo}

We optimize two complementary metrics: \emph{PR-AUC} (area under the
precision-recall curve, measuring document-level detection quality) and
\emph{Dice} ($2|P \cap G|/(|P|+|G|)$, measuring pixel-level mask overlap).
Our HPO formulation simultaneously maximizes:
\begin{equation}
  \max_{\boldsymbol{\theta} \in \Theta}\; \bigl(\text{PR-AUC}(\boldsymbol{\theta}),\;
  \text{Dice}(\boldsymbol{\theta})\bigr).
  \label{eq:moopt}
\end{equation}
Table~\ref{tab:hpo_space} defines the search space $\Theta$.

\begin{table}[!t]
  \centering
  \caption{HPO search space.}
  \label{tab:hpo_space}
  \begin{tabular}{llc}
    \toprule
    Parameter & Distribution & Range \\
    \midrule
    $\eta$ (learning rate) & log-uniform & $[5{\times}10^{-5},\; 9{\times}10^{-4}]$ \\
    $\lambda_{\text{wd}}$ (weight decay) & log-uniform & $[10^{-7},\; 10^{-4}]$ \\
    $p$ (dropout) & uniform & $[0.1,\; 0.4]$ \\
    $C_{\text{dec}}$ (decoder channels) & categorical & $\{96, 128, 192, 256\}$ \\
    $\lambda_{\text{mask}}$ (mask loss weight) & uniform & $[0.5,\; 3.0]$ \\
    \bottomrule
  \end{tabular}
\end{table}

We use MOTPE~\citep{motpe2022} as implemented in Optuna~\citep{optuna2019}.
After 50 trials on the inner splits, we select the configuration
$\boldsymbol{\theta}^*$ from the Pareto front $\mathcal{P}$ by minimizing the
Euclidean distance to the ideal point:
\begin{equation}
  \boldsymbol{\theta}^* = \argmin_{\boldsymbol{\theta} \in \mathcal{P}}
    \sqrt{(1 - \text{PR-AUC})^2 + (1 - \text{Dice})^2}.
  \label{eq:ideal_point}
\end{equation}
This criterion requires no user-specified weighting and assigns equal importance
to detection and localization. Figure~\ref{fig:pareto} illustrates the
resulting objective space.

\subsection{Nested Cross-Validation}
\label{subsec:ncv}

To obtain unbiased estimates while keeping HPO blind to test data, we use a
three-level nested CV scheme \citep{varma2006bias}. The outer loop uses 10
stratified folds; each fold serves once as the test set. For each outer fold,
the remaining data is split into 5 inner folds on which 50 MOTPE trials are
run (up to 15 epochs each, patience~$=1$). The optimal $\boldsymbol{\theta}^*$
is then used to retrain from scratch on all outer-fold training data (up to 100
epochs, early stopping patience~$=12$). This design ensures HPO never sees test data: 50 MOTPE trials per outer fold
yield 500 unique hyperparameter configurations in total (10 outer
folds $\times$ 50 trials); because each configuration is evaluated on 5 inner
splits, this corresponds to $10 \times 50 \times 5 = 2\,500$ individual model
trainings, plus 10 final models retrained on full outer-fold data.
